\documentclass[11pt]{article}
\usepackage[a4paper,margin=1in]{geometry}
\usepackage[T1]{fontenc}
\usepackage{lmodern,microtype}
\usepackage{amsmath,amssymb,amsthm,mathtools}
\usepackage{booktabs,enumitem}
\usepackage{xcolor}
\usepackage[numbers,sort&compress]{natbib}
\usepackage[colorlinks=true,linkcolor=blue!55!black,citecolor=blue!55!black]{hyperref}

\newtheorem{theorem}{Theorem}[section]
\newtheorem{proposition}[theorem]{Proposition}
\newtheorem{corollary}[theorem]{Corollary}
\theoremstyle{definition}
\newtheorem{definition}[theorem]{Definition}
\theoremstyle{remark}
\newtheorem{remark}[theorem]{Remark}

\title{Updated Anchored Head--Tail Certificate Ledger}
\author{Bin Wang}
\date{July 2026}

\begin{document}
\maketitle

\begin{abstract}
We update the exact finite-head/analytic-tail interface after four distinct
inputs.  RH-252 gives an all-order analytic existence interface for the
deterministic target, but its boundary constant $M_S$ is not numerically
certified.  RH-255 and RH-258 find no anchored head in their two expanded
admissibility classes at any of 32 endpoints.  RH-259 supplies a finite 23-endpoint
quotient block diagnostic whose worst root rate is $0.505642$, not a uniform
small-noise theorem.  We prove the updated gluing estimate and a scoped
zero-certificate corollary: the current audited classes have exactly zero
complete head--tail certificates.  This is a route ledger, not a global
nonexistence theorem.
\end{abstract}

\section{The updated certificate interface}

Let $\tau_{n,\theta}$ be the residual trace coefficients produced by a
candidate cloud or quotient at parameter $\theta$, and let $a_n$ be the
deterministic target coefficients of RH-243 \citep{WangRH243}.  Write
\begin{equation}
 L_\theta(z)=-\sum_{n\ge 2}\frac{\tau_{n,\theta}}{n}z^n,
 \qquad
 L_*(z)=-\sum_{n\ge 2}\frac{a_n}{n}z^n.
 \label{eq:logs}
\end{equation}
For a first omitted order $N\ge 3$, define the three nonnegative budgets
\begin{align}
 H_{<N}(\theta,R)&=\sum_{n=2}^{N-1}
     \frac{|\tau_{n,\theta}-a_n|R^n}{n},\\
 Q_{\ge N}(\theta,R)&=\sum_{n\ge N}
     \frac{|\tau_{n,\theta}|R^n}{n},\\
 A_{\ge N}(R)&=\sum_{n\ge N}
     \frac{|a_n|R^n}{n}.
 \label{eq:budgets}
\end{align}
The convention matters: an order-$2$--$12$ head has $N=13$.
This first-omitted-order formulation refines the bookkeeping interface of
RH-250 \citep{WangRH250}.

\begin{theorem}[Budgeted head--tail gluing]\label{thm:gluing}
Assume the series in \eqref{eq:logs} are absolutely convergent on
$|z|\le R$.  Then
\begin{equation}
 |L_\theta(z)-L_*(z)|
 \le H_{<N}(\theta,R)+Q_{\ge N}(\theta,R)+A_{\ge N}(R).
 \label{eq:logbound}
\end{equation}
If $H_{<N}\le h$, $Q_{\ge N}\le q$, and $A_{\ge N}\le t$ uniformly on a
parameter set, and if $|L_*(z)|\le B_*$ there, then with
$D_\theta=e^{L_\theta}$ and $D_*=e^{L_*}$,
\begin{equation}
 |D_\theta(z)-D_*(z)|
 \le e^{B_*}\bigl(e^{h+q+t}-1\bigr).
 \label{eq:detbound}
\end{equation}
Consequently a sequence of complete certificates yields locally uniform
determinant convergence whenever the three budgets tend uniformly to zero.
\end{theorem}

\begin{proof}
Subtract the two absolutely convergent series.  The terms with
$2\le n<N$ give the finite head $H_{<N}$; the remaining terms are
bounded separately by the absolute quotient and target tails.  This proves
\eqref{eq:logbound}.  For the determinant, put
$\Delta=L_\theta-L_*$.  Since $|L_*|\le B_*$,
\[
 |e^{L_\theta}-e^{L_*}|
 =|e^{L_*}|\,|e^\Delta-1|
 \le e^{B_*}(e^{|\Delta|}-1),
\]
and \eqref{eq:detbound} follows from \eqref{eq:logbound}.
\end{proof}

\begin{proposition}[Analytic target tail exists, but its constant is open]
\label{prop:target}
RH-252 \citep{WangRH252} gives a scaled zero-free radius
\[
 \rho_*=r_H\lambda=1.42678748386407\ldots>1,
 \qquad r_H=0.85.
\]
For every $0\le R<S<\rho_*$, the normalized logarithm has a finite
boundary supremum $M_S$ and
\begin{equation}
 A_{\ge N}(R)\le M_S\frac{(R/S)^N}{1-R/S}.
 \label{eq:targettail}
\end{equation}
At $R=1$, $N=13$, and $S=1.35$, the geometric factor per $M_S$ is
$0.07796985628233917$.  The existence of $M_S$ is exact; no certified
numerical upper bound for it is available in the archived audit.
\end{proposition}

\begin{proof}
This is the Cauchy estimate for the holomorphic logarithm on the disk supplied
by RH-252.  The coefficient of $z^n$ is $-a_n/n$, so summing the resulting
geometric series gives \eqref{eq:targettail}.  The final numerical factor is
only the displayed geometric factor; multiplying it by a truncated boundary
scan would not certify $M_S$.
\end{proof}

\section{Scoped zero-certificate result}

The head classes are deliberately kept separate.  RH-255 uses a convex
single-use shell box containing every prefix and binary shell subset.  Its
archived primal--dual audit has 0 passes at all 32 endpoints.  RH-258 tests the
unit-cap signed-integer lattice $\{-1,0,1\}^J$, with exact-integrality MILP
solutions and zero reported MIP gap; it also has 0 passes at all 32 endpoints.
Neither audit supplies an operator realization for an arbitrary mask
\citep{WangRH255,WangRH258}.

\begin{corollary}[Zero complete certificates in the audited classes]
\label{cor:zero}
Restrict the candidate head to the union of the RH-255 box class and the
RH-258 unit-cap signed-integer class, with the archived 32-endpoint tolerance
protocol.  There is no complete head--tail certificate in this restricted
class.
\end{corollary}

\begin{proof}
Every complete certificate has, as a necessary first component, a passing
anchored finite head.  The RH-255 audit supplies 0/32 such passes and the
RH-258 audit supplies 0/32.  Therefore their union supplies no passing head,
so the set of complete certificates is empty, independently of the tail
budgets.  This conclusion is scoped to the two archived classes and their
finite tolerance protocol.
\end{proof}

The two classes cover 64 class--endpoint cases.  The RH-255 audit excludes
\[
 62{,}030{,}604{,}700
\]
eligible binary subsets through the box obstruction.  RH-258's unit-cap
lattice contains
\[
 39{,}417{,}456{,}084{,}975{,}216
\]
signed points in aggregate.  These counts are finite audit descriptors, not
claims about all possible selectors.

\section{Quotient-tail ledger}

RH-259 extends the ordered-Schur quotient calculation to dimension 1024
\citep{WangRH246,WangRH259}.  All
23 audited twelfth powers are contractive, but the worst root rate is
\[
 q_{12}=0.5056418005507071,
\]
at the right channel with $\sigma=0.0025$ and dimension 1024.  The finite
unit-disk logarithmic tail diagnostic is
$5.654507945432548\times10^{-4}$.  Nine archived endpoints remain outside
the calculation, and the floating matrix norms are not interval enclosures.
Thus this is a family of finite diagnostics, not a uniform small-noise
quotient certificate.

\begin{center}
\small
\begin{tabular}{p{0.46\linewidth}p{0.39\linewidth}}
\toprule
certificate component & archived status \\
\midrule
deterministic target analytic tail & exact existence; $M_S$ not certified \\
RH-255 expanded single-use head & 0/32 passes \\
RH-258 unit-cap signed-integer head & 0/32 passes \\
RH-259 quotient block tail & 23 finite endpoints; uniformity open \\
cloud coefficient bridge & open \\
complete head--tail certificate & 0 \\
\bottomrule
\end{tabular}
\end{center}

\section{Protocol, boundary, and next input}

The experiment reads the four archived JSON ledgers, checks 44 source
consistency predicates, and writes a new structured ledger.  It does not
rerun the dimension-1024 Schur calculation or fit new moments.  The result
therefore separates exact inherited interfaces from finite numerical audits:
the target-tail existence flag is true, while a certified boundary constant,
a legal head pass, a cloud coefficient bridge, and a uniform quotient tail
are all false.

The obstruction does not exclude larger integer caps, signed or complex
selectors derived from an invariant quotient, larger resolved windows, a
future interval bound for $M_S$, or another operator realization.  It also
does not establish global nonexistence.  Gates A--E remain false/open.  No
Hilbert--Polya operator is constructed; no Riemann zeros are identified; no
equality with the completed zeta divisor and no implication of RH is claimed.

The next reopening input must supply at least one genuinely new legal head,
a certified $M_S$, a uniform quotient theorem, or an exact coefficient bridge;
the useful route requires the relevant obligations together rather than a
finite-order extrapolation.

\bibliographystyle{plainnat}
\bibliography{references}
\end{document}
